
    %%%%%%%%%%%%%%%%%%%%%%%%%%%%%%%%%%%%%%%%%%%%%%%%%%%%%%%%%%%%%%%%%%%%%%%%%%%%%%%%
    %%% Classe do documento
    \documentclass[12pt, addpoints]{exam}
    \UseRawInputEncoding

    %%%%%%%%%%%%%%%%%%%%%%%%%%%%%%%%%%%%%%%%%%%%%%%%%%%%%%%%%%%%%%%%%%%%%%%%%%%%%%%%
    % preambulo.tex
    %
    % Arquivo de configuração de packages para uso o LaTeX, conforme minhas
    % preferências e modelos pessoais.
    %
    % Para maiores informações, visite:
    %    https://github.com/abrantesasf/latex
    %
    % NÃO ALTERE SE NÃO SOUBER O QUE ESTÁ FAZENDO!
    %%%%%%%%%%%%%%%%%%%%%%%%%%%%%%%%%%%%%%%%%%%%%%%%%%%%%%%%%%%%%%%%%%%%%%%%%%%%%%%%


    %%%%%%%%%%%%%%%%%%%%%%%%%%%%%%%%%%%%%%%%%%%%%%%%%%%%%%%%%%%%%%%%%%%%%%%%%%%%%%%%
    %%% Estruturas de controle:
    \input{static/utils/controles}

    %%%%%%%%%%%%%%%%%%%%%%%%%%%%%%%%%%%%%%%%%%%%%%%%%%%%%%%%%%%%%%%%%%%%%%%%%%%%%%%%
    %%% Compilação condicional em PDF:
    \input{static/utils/pdf}

    %%%%%%%%%%%%%%%%%%%%%%%%%%%%%%%%%%%%%%%%%%%%%%%%%%%%%%%%%%%%%%%%%%%%%%%%%%%%%%%%
    %%% Configurações de layout da página:
    \input{static/utils/layout}

    %%%%%%%%%%%%%%%%%%%%%%%%%%%%%%%%%%%%%%%%%%%%%%%%%%%%%%%%%%%%%%%%%%%%%%%%%%%%%%%%
    %%% Configurações para autores:
    \input{static/utils/autores}

    %%%%%%%%%%%%%%%%%%%%%%%%%%%%%%%%%%%%%%%%%%%%%%%%%%%%%%%%%%%%%%%%%%%%%%%%%%%%%%%%
    %%% Cabeçalho e rodapé:
    %%% Atenção! É MUITO DIFÍCIL ajustar apropriadamente os running headers
    %%% e running footers da primeira vez. Você pode confiar no LaTeX e deixar que
    %%% ele faça o ajuste automaticamente ou pode quebrar a cabeça e
    %%% tentar melhorar algo que o LaTeX já faz muito bem, mesmo que não fique
    %%% exatamente do jeito que você quer. O que você prefere? Se quiser ajustar
    %%% manualmente, descomente a linha abaixo e faça as configurações no arquivo
    %%% "utils/headings.tex".
    %\input{static/utils/headings}

    %%%%%%%%%%%%%%%%%%%%%%%%%%%%%%%%%%%%%%%%%%%%%%%%%%%%%%%%%%%%%%%%%%%%%%%%%%%%%%%%
    %%% Configuração para as fontes do documento:
    %%% Se você quiser usar fontes próprias ou do sistema, precisará ajustar as
    %%% configurações no arquivo "utils/fontes.tex".
    %%% ATENÇÃO: por padrão eu utilizo XeLaTeX com fontes proprietárias projetadas
    %%% por Matthew Butterick, adquiridas em https://mbtype.com, que não podem ser
    %%% distribuídas devido à licença de utilização. Se você usar XeLaTeX, você
    %%% DEVERÁ OBRIGATORIAMENTE ajustar as fontes no arquivo "utils/fontes.tex".
    \input{static/utils/fontes}

    %%%%%%%%%%%%%%%%%%%%%%%%%%%%%%%%%%%%%%%%%%%%%%%%%%%%%%%%%%%%%%%%%%%%%%%%%%%%%%%%
    %%% Sumário:
    \input{static/utils/toc}

    %%%%%%%%%%%%%%%%%%%%%%%%%%%%%%%%%%%%%%%%%%%%%%%%%%%%%%%%%%%%%%%%%%%%%%%%%%%%%%%%
    %%% Matemática:
    \input{static/utils/matematica}

    %%%%%%%%%%%%%%%%%%%%%%%%%%%%%%%%%%%%%%%%%%%%%%%%%%%%%%%%%%%%%%%%%%%%%%%%%%%%%%%%
    %%% Notas de rodapé:
    \input{static/utils/footnotes}

    %%%%%%%%%%%%%%%%%%%%%%%%%%%%%%%%%%%%%%%%%%%%%%%%%%%%%%%%%%%%%%%%%%%%%%%%%%%%%%%%
    %%% Referências cruzadas, links, citações:
    \input{static/utils/crossrefs}

    %%%%%%%%%%%%%%%%%%%%%%%%%%%%%%%%%%%%%%%%%%%%%%%%%%%%%%%%%%%%%%%%%%%%%%%%%%%%%%%%
    %%% Configurações para os hiperlinks em PDF:
    \input{static/utils/pdfconfig}

    %%%%%%%%%%%%%%%%%%%%%%%%%%%%%%%%%%%%%%%%%%%%%%%%%%%%%%%%%%%%%%%%%%%%%%%%%%%%%%%%
    %%% Glossário e índice remissivo:
    \input{static/utils/glossindex}

    %%%%%%%%%%%%%%%%%%%%%%%%%%%%%%%%%%%%%%%%%%%%%%%%%%%%%%%%%%%%%%%%%%%%%%%%%%%%%%%%
    %%% Referências bibliográficas:
    \input{static/utils/bibliografia}

    %%%%%%%%%%%%%%%%%%%%%%%%%%%%%%%%%%%%%%%%%%%%%%%%%%%%%%%%%%%%%%%%%%%%%%%%%%%%%%%%
    %%% Cores:
    \input{static/utils/cores}

    %%%%%%%%%%%%%%%%%%%%%%%%%%%%%%%%%%%%%%%%%%%%%%%%%%%%%%%%%%%%%%%%%%%%%%%%%%%%%%%%
    %%% Computação:
    \input{static/utils/computacao}

    %%%%%%%%%%%%%%%%%%%%%%%%%%%%%%%%%%%%%%%%%%%%%%%%%%%%%%%%%%%%%%%%%%%%%%%%%%%%%%%%
    %%% Gráficos:
    \input{static/utils/graficos}

    %%%%%%%%%%%%%%%%%%%%%%%%%%%%%%%%%%%%%%%%%%%%%%%%%%%%%%%%%%%%%%%%%%%%%%%%%%%%%%%%
    %%% Ícones:
    \input{static/utils/icones}

    %%%%%%%%%%%%%%%%%%%%%%%%%%%%%%%%%%%%%%%%%%%%%%%%%%%%%%%%%%%%%%%%%%%%%%%%%%%%%%%%
    %%% Blocos:
    \input{static/utils/blocos}

    %%%%%%%%%%%%%%%%%%%%%%%%%%%%%%%%%%%%%%%%%%%%%%%%%%%%%%%%%%%%%%%%%%%%%%%%%%%%%%%%
    %%% Boxes coloridos:
    \input{static/utils/colorbox}

    %%%%%%%%%%%%%%%%%%%%%%%%%%%%%%%%%%%%%%%%%%%%%%%%%%%%%%%%%%%%%%%%%%%%%%%%%%%%%%%%
    %%% Tabelas:
    \input{static/utils/tabelas}

    %%%%%%%%%%%%%%%%%%%%%%%%%%%%%%%%%%%%%%%%%%%%%%%%%%%%%%%%%%%%%%%%%%%%%%%%%%%%%%%%
    %%% Ambientes de listas:
    \input{static/utils/listas}

    %%%%%%%%%%%%%%%%%%%%%%%%%%%%%%%%%%%%%%%%%%%%%%%%%%%%%%%%%%%%%%%%%%%%%%%%%%%%%%%%
    %%% Ambientes floats e similares:
    \input{static/utils/floats}

    %%%%%%%%%%%%%%%%%%%%%%%%%%%%%%%%%%%%%%%%%%%%%%%%%%%%%%%%%%%%%%%%%%%%%%%%%%%%%%%%
    %%% Ajustes para a classe EXAM:
    \input{static/utils/exam}

    %%%%%%%%%%%%%%%%%%%%%%%%%%%%%%%%%%%%%%%%%%%%%%%%%%%%%%%%%%%%%%%%%%%%%%%%%%%%%%%%
    %%% Ajustes para textos:
    \input{static/utils/texto}

    %%%%%%%%%%%%%%%%%%%%%%%%%%%%%%%%%%%%%%%%%%%%%%%%%%%%%%%%%%%%%%%%%%%%%%%%%%%%%%%%
    %%% Ambientes, aliases, comandos e customizações em geral para serem
    %%% utilizadas nos documentos. Defina aqui qualquer customização específica
    %%% para seus documentos!
    \input{static/utils/customizacoes}

    %%%%%%%%%%%%%%%%%%%%%%%%%%%%%%%%%%%%%%%%%%%%%%%%%%%%%%%%%%%%%%%%%%%%%%%%%%%%%%%%
    %%% Ajuste do layout, espaçamento de linhas e etc.:
    \geometry{a4paper, portrait, left=2.0cm, right=2.0cm, top=2.0cm, bottom=2.0cm}
    \singlespacing

    %%%%%%%%%%%%%%%%%%%%%%%%%%%%%%%%%%%%%%%%%%%%%%%%%%%%%%%%%%%%%%%%%%%%%%%%%%%%%%%%
    %%% Configurações de cabeçalho e rodapé (não use fancyhdr pois ocorre conflito):
    \pagestyle{headandfoot}
    \runningheadrule
    \firstpageheader{}{}{}
    \runningheader{Sem disciplina}{}{Avaliação}
    \firstpagefooter{}{}{Página \thepage\ de \numpages}
    \runningfooter{}{}{Página \thepage\ de \numpages}
    \runningfootrule

    %%%%%%%%%%%%%%%%%%%%%%%%%%%%%%%%%%%%%%%%%%%%%%%%%%%%%%%%%%%%%%%%%%%%%%%%%%%%%%%%
    %%% Configurações para as propriedades do PDF:
    \hypersetup{
    hidelinks,           % Comente para web, descomente para impressão
    colorlinks=true,      % True para web, False para impressão
    pdftitle=VaiQueDa: Avaliação,
    pdfauthor=Matheus Endlich Silveira,
    pdfsubject=Sem disciplina,
    pdfkeywords=['Sem tag'],
    pdfinfo={
        CreationDate={}, % Ex.: D:AAAAMMDDHH24MISS
        ModDate={}       % Ex.: D:AAAAMMDDHH24MISS
    }
    }
    \input{static/utils/pdfconfig}

    %%%%%%%%%%%%%%%%%%%%%%%%%%%%%%%%%%%%%%%%%%%%%%%%%%%%%%%%%%%%%%%%%%%%%%%%%%%%%%%%
    %%% Começa o documento
    \begin{document}

    %%%%%%%%%%%%%%%%%%%%%%%%%%%%%%%%%%%%%%%%%%%%%%%%%%%%%%%%%%%%%%%%%%%%%%%%%%%%%%%%
    %%% Folha de instruções padronizadas da UVV para a prova
    \pdfbookmark[1]{Instruções}{instrucoes}
    \begin{figure}[H]
        \begin{center}
            \includegraphics[scale=0.3]{{static/imgs/logo-uvv.png}}
        \end{center}	
    \end{figure}

    \vspace{-1cm}

    \begin{table}[H]
        \centering
        \begin{tabular}{|llll|}
            \hline
            \multicolumn{2}{|l|}{\textbf{Disciplina:} Sem disciplina} & 
            \multicolumn{1}{l|}{\multirow{3}{*}{\textbf{\makecell{Nota:\\ \,\\ \,}}}} & 
            \multirow{3}{*}{\textbf{\makecell{Coordenador:\\ \,\\ \,}}} \\ 
            \cline{1-2}
            \multicolumn{2}{|l|}{\textbf{Professor:} Matheus Endlich Silveira \hspace{4.72cm} } & 
            \multicolumn{1}{l|}{} & \\ 
            \cline{1-2}
            \multicolumn{2}{|l|}{\textbf{Aluno:}} & 
            \multicolumn{1}{l|}{} & \\ 
            \hline
            \multicolumn{1}{|l|}{\textbf{Turma:} \hspace{6.3cm} } & 
            \multicolumn{1}{l|}{\textbf{Semestre:}} & 
            \multicolumn{2}{l|}{\textbf{Valor:} 10 pontos} \\ 
            \hline
            \multicolumn{1}{|l|}{\textbf{Data:}} & 
            \multicolumn{3}{l|}{\textbf{Avaliação:}} \\ 
            \hline
        \end{tabular}
    \end{table}

    \begin{center}
        \fbox{\setlength\fboxsep{0.3cm} \fbox{\parbox{15.4cm}{
            \begin{center}
                \textbf{INSTRUÇÕES PARA A REALIZAÇÃO DESTA AVALIAÇÃO:}
            \end{center}
            \begin{itemize}[leftmargin=*]
                \item Esta é a primeira avaliação da disciplina Sem disciplina, 
                    e engloba o conteúdo referente Sem tag.
                \item Esta avaliação contém 20 questões objetivas, \textbf{todas obrigatórias}.
                \item \textbf{Leia atentamente cada questão!} As questões objetivas terão uma,
                    e apenas uma única, resposta a ser marcada.
                \item \textbf{Desligue o celular} antes de começar. A avaliação é
                    \textbf{individual} e \textbf{sem consulta}.
                \item As questões podem ser respondidas, na prova, com lápis ou caneta. Ao final
                    da prova \textbf{{VOCÊ DEVERÁ PREENCHER O GABARITO COM CANETA
                    DE COR PRETA}} \textbf{\underline{OU AZUL}} para que seja feita a correção
                    automatizada através da leitura do padrão de respostas marcado no
                    gabarito.
                \item \textbf{CUIDADO ao preencher o gabarito} pois eles são \textbf{INDIVIDUAIS
                    e NOMEADOS} (cada estudante tem um gabarito próprio específico, com um
                    código de identificação único). \textbf{Questões rasuradas são anuladas}
                    pelo software de leitura do gabarito.
                \item Ao final da prova, \textbf{DEVOLVA A PROVA E O GABARITO} para o professor.
                \item Siga todas as normas de \textbf{Integridade Acadêmica} da disciplina.
                    Alunos que forem flagrados com qualquer espécie de ``cola'' ou trocando
                    informações com outros alunos terão suas avaliações recolhidas, as notas
                    zeradas, e a situação será encaminhada para a coordenação para a aplicação
                    das penalidades previstas pela universidade.
                \item Com 90 minutos de avaliação você terá tempo de até 4,min por questão,
                    em média.
                \item Boa avaliação!
            \end{itemize}
            \vspace{2cm}
        }}}
    \end{center}
    \newpage

    %%%%%%%%%%%%%%%%%%%%%%%%%%%%%%%%%%%%%%%%%%%%%%%%%%%%%%%%%%%%%%%%%%%%%%%%%%%%%%%%
    %%% Inicia o bloco de questões:
    \begin{questions}

    \newpage
    \question

A interposição de um circuito de memória cache entre o processador e a memória principal (RAM)
\begin{choices}
\choice diminui o tráfego de instruções e/ou dados no barramento de memória
\choice aumenta o tráfego de instruções e/ou dados entre memória e disco
\choice diminui o tráfego de instruções e/ou dados entre memória e disco
\choice permite acessos concorrentes à memória RAM
\choice aumenta o tráfego de instruções e/ou dados no barramento de memória
\end{choices}

\question

Qual das seguintes afirmações sobre crescimento assintótico de funções não é verdadeira:
\begin{choices}
\choice Se \( f(n) = O(g(n)) \) e \( g(n) = O(h(n)) \) então \( f(n) = O(h(n)) \)
\choice \( \log n^2 = O(\log n) \)
\choice Se \( f(n) = O(g(n)) \) então \( g(n) = O(f(n)) \)
\choice \( 2^{n+1} = O(2^n) \)
\choice \( 2n^2 + 3n + 1 = O(n^2) \)
\end{choices}

\question

Considere um problema em que são dados 5 objetos com os seguintes pesos e valores:\\

\text{pesos: } $(W_1, W_2, W_3, W_4, W_5) = (6, 10, 9, 5, 12)$\\

\text{valores: } $(P_1, P_2, P_3, P_4, P_5) = (8, 5, 10, 15, 7)$\\

Além disso, é dada uma mochila que suporta até 30 unidades de peso, para transportar os objetos. O objetivo do problema é preencher a mochila de tal forma que o valor total dos objetos a serem transportados seja o maior possível, mas sem exceder o limite de peso suportado pela mochila. Assuma que é permitido colocar fração de um objeto na mochila. Qual das seguintes alternativas corresponde ao valor máximo obtido no preenchimento da mochila:
\begin{choices}
\choice 12.2
\choice 21.5
\choice 43.1
\choice 38.83
\choice 30.34
\end{choices}

\question

Quais dos algoritmos de ordenação abaixo possuem tempo no pior caso e tempo médio de execução proporcional a \( O(n \log n) \).
\begin{choices}
\choice Heap sort e selection sort
\choice Merge sort e heap sort
\choice Quicksort e merge sort
\choice Merge sort e bubble sort
\choice Bubble sort e Quick sort
\end{choices}

\question

O processo de visualização de objetos 3D envolve uma série de passos desde a representação vetorial de um objeto até a exibição da imagem correspondente na tela do computador (pipeline 3D). Selecione a alternativa abaixo que reflete a ordem correta em que esses passos devem ocorrer. 
\begin{choices}
\choice Transformação de câmera, mapeamento para coordenadas de tela, recorte 3D, rasterização, projeção. 
\choice Transformação de câmera, recorte 3D, projeção, mapeamento para coordenadas de tela, rasterização. 
\choice Recorte 3D, transformação de câmera, rasterização, projeção, mapeamento para coordenadas de tela. 
\choice Projeção, transformação de câmera, recorte 3D, mapeamento para coordenadas de tela, rasterização. 
\choice Nenhuma das respostas acima está correta.
\end{choices}

\question

Seja a árvore binária abaixo a representação de um espaço de estados para um problema \( p \), em que o estado inicial é \( a \), e \( i \) e \( f \) são estados finais.



\begin{figure}[H]

    \centering

    \includegraphics[width=0.5\linewidth]{static/imgs/defaul.png}

    \caption{Questão 59}

    \label{fig:enter-label}

\end{figure}



Um algoritmo de busca em largura-primeiro forneceria a seguinte sequência de estados como primeira alternativa a um caminho-solução para o problema \( p \):
\begin{choices}
\choice \( a \ b \ c \ d \ e \ f \)
\choice \( a \ b \ d \ h \ e \ i \)
\choice \( a \ c \ f \)
\choice \( a \ b \ e \ i \)
\choice \( a \ b \ d \ e \ f \)
\end{choices}

\question

Na linguagem SQL, qual das seguintes palavras-chave é usada para modificar os

dados em uma tabela existente?


\begin{choices}
\choice CREATE
\choice INSERT
\choice UPDATE
\choice Nenhuma das alternativas acima
\choice SELECT
\end{choices}

\question

\textbf Supondo a Relação \texttt{PROJ (PNO, Nome, Orçam)}, com chave primária \texttt{PNO} e a Relação \texttt{DSG (ENO, PNO, Dur, Resp)}, com chave primária \{\texttt{ENO, PNO}\} e chave estrangeira \texttt{PNO} em relação a \texttt{PROJ}, a asserção abaixo \textbf{NÃO} expressa:



\begin{equation}\forall g \in \texttt{DSG}, \exists j \in \texttt{PROJ} : g.\texttt{PNO} = j.\texttt{PNO}\end{equation}
\begin{choices}
\choice Uma restrição de integridade de chave estrangeira em $\texttt{DSG}$.
\choice Uma restrição a ser verificada na atualização de tuplas em $\texttt{DSG}$.
\choice Uma restrição de integridade de chave primária em $\texttt{PROJ}$.
\choice Uma restrição que define um estado consistente do banco de dados.
\choice Uma restrição a ser verificada na inserção de tuplas em $\texttt{DSG}$.
\end{choices}

\question

O contador da figura abaixo é:

\begin{figure}

    \centering

    \includegraphics[width=0.5\linewidth]{static/imgs/defaul.png}

    \label{fig:enter-label}

\end{figure}
\begin{choices}
\choice assíncrono
\choice anisócrono
\choice síncrono
\choice isócrono
\choice auto-sincronizado
\end{choices}

\question

Considere as seguintes tabelas em uma base de dados relacional:\\

Departamento (\underline{CodDepto}, NomeDepto)\\

Empregado (\underline{CodEmp}, NomeEmp, CodDepto, Salario)\\

Considere a seguinte consulta SQL:

\begin{verbatim}

SELECT d.CodDepto, d.NomeDepto, SUM(e.Salario)

FROM Departamento d

JOIN Empregado e ON d.CodDepto = e.CodDepto

GROUP BY d.CodDepto, d.NomeDepto

HAVING COUNT(e.CodEmp) > 2 AND AVG(e.Salario) > 40;

\end{verbatim}



Qual o resultado da consulta?
\begin{choices}
\choice Para cada departamento que tem mais que dois empregados e cuja média salarial é maior que 40, obter um grupo de linhas que contém, para cada empregado do departamento, o código de seu departamento, seguido do nome de seu departamento, seguido da soma dos salários dos empregados do departamento.
\choice Para cada departamento que tem mais que dois empregados e cuja média salarial é maior que 40, obter o código de departamento, seguido do nome do departamento, seguido da soma dos salários dos empregados do departamento.
\choice A consulta não retorna nada pois está incorreta.
\choice Para cada departamento que tem mais que dois empregados e cuja média salarial, considerando todos empregados do departamento, exceto os dois primeiros, é maior que 40, obter o código de departamento, seguido do nome do departamento, seguido da soma dos salários dos empregados do departamento.
\choice Para cada empregado que tem mais que dois departamentos, ambos com média salarial maior que 40, obter o código de departamento, seguido do nome do departamento, seguido da soma dos salários dos empregados do departamento.
\end{choices}

\question

Considere as seguintes afirmativas:

\begin{itemize}

    \item O modelo matemático de uma lista é a sequência linear de itens, cuja principal propriedade estrutural é a posição relativa dos elementos dentro da sequência.

    \item A fila e a pilha são consideradas casos especiais da lista.

    \item Numa fila a inserção e a retirada são feitas no mesmo extremo.

    \item Numa lista a inserção e a retirada podem ser feitas em qualquer posição.

    \item Numa pilha apenas a inserção pode ser feita em qualquer posição.

\end{itemize}



Quais são as afirmativas verdadeiras?
\begin{choices}
\choice somente II, III e IV
\choice somente II, IV e V
\choice todas
\choice somente I e III
\choice somente I, II e IV
\end{choices}

\question

Professor Mac Sperto propôs o seguinte algoritmo de ordenação, chamado de Super Merge, similar ao \textit{merge sort}: divida o vetor em 4 partes do mesmo tamanho (ao invés de 2, como é feito no \textit{merge sort}). Ordene recursivamente cada uma das partes e depois intercale-as por um procedimento semelhante ao procedimento de intercalação do \textit{merge sort}. Qual das alternativas abaixo é verdadeira?
\begin{choices}
\choice Nenhuma das afirmações acima está correta
\choice Super Merge está correto, mas consome tempo menor que \( O(\text{merge sort}) \)
\choice Super Merge está correto, mas consome tempo \( O(\text{merge sort}) \)
\choice Super Merge está correto, mas consome tempo maior que \( O(\text{merge sort}) \)
\choice Super Merge não está correto. Não é possível ordenar quebrando o vetor em 4 partes
\end{choices}

\question

Se \( \lim_{n \to \infty} \left( \frac{1}{n^2} + \frac{2}{n^2} + \cdots + \frac{n-1}{n^2} \right) = L \) então


\begin{choices}
\choice \( L = 2 \)
\choice \( L = 1/2 \)
\choice \( L = \infty \)
\choice \( L = 1 \)
\choice \( L = 0 \)
\end{choices}

\question

Quais das seguintes afirmações são verdadeiras? As Métricas de software servem para:



\begin{itemize}

    \item[I)] indicar a qualidade do produto e avaliar a produtividade.

    \item[II)] auxiliar na melhoria do processo.

    \item[III)] formar uma base para as estimativas e justificar a aquisição de ferramentas.

    \item[IV)] determinar se a utilização de um método traz benefícios ou não.

\end{itemize}
\begin{choices}
\choice Todas as alternativas.
\choice Apenas as alternativas I, IV.
\choice Apenas as alternativas II e III.
\choice Apenas as alternativas I, II e IV.
\choice Nenhuma delas.
\end{choices}

\question

Ao analisar um novo sistema de gerenciamento de bancos de dados que surgiu no

mercado, você percebe que a entrada às operações do modelo de dados são tabelas

mas as saídas nunca são tabelas, são sempre outras estruturas. O vendedor desse

SGBD afirma que ele é um SGBD relacional. Pode-se afirmar que:


\begin{choices}
\choice Nenhuma das alternativas anteriores.
\choice Esse SGBD não é relacional pois, para ser verdadeiramente relacional, as operações do modelo de dados não podem trabalhar recebendo tabelas como entrada.
\choice Esse SGBD é relacional pois consegue trabalhar com tabelas para a entrada das operações do modelo de dados. A saída não é importante.
\choice Esse SGBD não é relacional pois no modelo relacional a única estrutura percebida pelos usuários são tabelas.
\choice Esse SGBD é relacional pois o vendedor afirmou que é.
\end{choices}

\question

Considere \( C(x) \) uma função que define a complexidade de um problema \( x \); \( E(x) \) uma função que define o esforço (em termos de tempo) exigido para se resolver o problema \( x \). Sejam dois problemas denominados \( p1 \) e \( p2 \). Assinale a alternativa correta.


\begin{choices}
\choice \( C(p1 + p2) < C(p1) + C(p2) \)
\choice Se \( C(p1) < C(p2) \) então \( E(p1) > E(p2) \)
\choice Se \( C(p1) < C(p2) \) então \( E(p1) < E(p2) \)
\choice Nenhuma das alternativas anteriores
\choice \( E(p1 + p2) < E(p1) + E(p2) \)
\end{choices}

\question

Filtro da mediana é:
\begin{choices}
\choice Nenhuma das respostas acima.
\choice Indicado para atenuar ruído com preservação de bordas (i.e., rápidas transições de nível em imagens).
\choice Indicado para detectar tonalidades específicas em uma imagem.
\choice Indicado para detectar bordas em imagens.
\choice Indicado para detectar formas específicas em imagens.
\end{choices}

\question

Seja a seguinte linguagem, onde \(\epsilon\) representa o string vazio e \$ representa um marcador de fim de entrada:



$S \rightarrow ABCD \\$\\

$A \rightarrow a \mid \epsilon \\$\\

$B \rightarrow a \mid \epsilon \\$\\

$C \rightarrow c \mid \epsilon \\$\\

$D \rightarrow S \mid c \mid \epsilon$\\



É incorreto afirmar que:
\begin{choices}
\choice O conjunto FOLLOW(D) é igual a FOLLOW(S)
\choice O conjunto FOLLOW(A) = a, c, \$
\choice O conjunto FIRST(A) = a, \(\epsilon\)
\choice O conjunto FIRST(D) é igual ao conjunto FIRST(S)
\choice O conjunto FOLLOW(B) = c, \$
\end{choices}

\question

Os protocolos de transporte atribuem a cada serviço um identificador único, o qual é empregado para encaminhar uma requisição de um aplicativo cliente ao processo servidor correto. Nos protocolos de transporte TCP e UDP, como esse identificador se denomina?
\begin{choices}
\choice Endereço IP
\choice Porta
\choice Protocolo de aplicação
\choice Identificador do processo (PID)
\choice Conexão
\end{choices}

\question

Considere as seguintes tabelas em uma base de dados relacional:\\

\textbf{Departamento (CodDepto, NomeDepto)}\\

\textbf{Empregado (CodEmp, NomeEmp, CodDepto)}\\

Deseja-se obter uma tabela na qual cada linha é a concatenação de uma linha da tabela Departamento com uma linha da tabela de Empregado. Caso um departamento não possua empregados, sua linha no resultado deve conter vazio (NULL) nos campos referentes ao empregado. A operação de álgebra relacional que deve ser aplicada para combinar estas duas tabelas é:
\begin{choices}
\choice Projeção
\choice Divisão
\choice União
\choice Junção interna
\choice Junção externa
\end{choices}



    %%%%%%%%%%%%%%%%%%%%%%%%%%%%%%%%%%%%%%%%%%%%%%%%%%%%%%%%%%%%%%%%%%%%%%%%%%%%%%%%
    %%% Finaliza o bloco de questões:
    \end{questions}
    \end{document}
    